\documentclass[11pt]{article}
\usepackage[margin=1in]{geometry}
\usepackage{amsmath, amssymb}
\usepackage{booktabs}
\usepackage{hyperref}
\usepackage{enumitem}
\usepackage{setspace}
\setstretch{1.15}
\setlength{\parindent}{0pt}
\setlist[itemize]{itemsep=0.25em, topsep=0.25em}

\begin{document}

{\Large \textbf{ECON 204 (Contemporary Economic Issues): In-Class Empirical Group Activity}}\\
{\large \textbf{Inflation, Cost of Living, and Monetary Policy in Canada}}\\[0.5em]
\textbf{Format:} groups of 3--4 \hspace{1.5em} \textbf{Time:} 45--60 minutes \hspace{1.5em} \textbf{Tool:} Excel or Google Sheets

\vspace{0.75em}
\textbf{Learning goals.} In this activity you will (i) construct Canadian inflation from CPI data, (ii) connect inflation to cost-of-living changes, and (iii) evaluate how the Bank of Canada’s policy rate moves relative to inflation.

\vspace{0.75em}
\textbf{Data you will use (provided by the instructor).} You will receive four files:
\begin{itemize}
\item \texttt{canada\_cpi\_allitems\_monthly\_2015\_2025.csv}: all-items CPI level for Canada (monthly).
\item \texttt{canada\_policy\_rate\_target\_monthly\_eom\_2015\_2025.csv}: policy interest rate (``target rate''), end-of-month (monthly).
\item \texttt{canada\_avg\_hourly\_earnings\_hourly\_paid\_monthly\_2015\_2025.csv}: a simple wage proxy (monthly).
\item \texttt{canada\_inflation\_policy\_wage\_monthly\_2015\_2025.csv}: a cleaned file that merges the three series by month.
\end{itemize}

\section*{Part A. Organize the dataset (10 minutes)}
Open the Excel template provided by the instructor. Confirm that the \texttt{Data} sheet contains these columns: \texttt{date}, \texttt{cpi}, \texttt{policy\_rate}, \texttt{avg\_hourly\_earnings}. 

\vspace{0.5em}
\textbf{Checkpoint A1.} What is the first month in the dataset? What is the last month?\\[0.75em]
\underline{\hspace{0.9\linewidth}}\\[0.75em]
\underline{\hspace{0.9\linewidth}}

\section*{Part B. Construct inflation (10--15 minutes)}
Use the CPI series to compute year-over-year inflation:
\[
\pi_t = \left(\frac{CPI_t}{CPI_{t-12}} - 1\right)\times 100.
\]
In the template, this is computed on the \texttt{Calculations} sheet.

\vspace{0.5em}
\textbf{Checkpoint B1.} Identify the month with the highest year-over-year inflation in 2015--2025. Record the month and the inflation rate.\\[0.75em]
\underline{\hspace{0.9\linewidth}}\\[0.75em]
\underline{\hspace{0.9\linewidth}}

\section*{Part C. Make two charts (10--15 minutes)}
Create the following two charts (titles and axis labels matter):
\begin{itemize}
\item \textbf{Chart 1:} CPI level and inflation (YoY). If you use a secondary axis, explain what each axis measures.
\item \textbf{Chart 2:} Inflation (YoY) and the policy rate. Use the same monthly date axis.
\end{itemize}

\vspace{0.5em}
\textbf{Checkpoint C1.} In Chart 2, does the policy rate appear to rise \emph{before} inflation peaks, \emph{after} inflation peaks, or roughly \emph{at the same time}? Give one sentence.\\[0.75em]
\underline{\hspace{0.9\linewidth}}\\[0.75em]
\underline{\hspace{0.9\linewidth}}

\section*{Part D. Cost-of-living interpretation (10 minutes)}
Inflation is an economy-wide average, but cost-of-living pressure can be felt very differently across households. Use your charts to anchor your answers.

\vspace{0.5em}
\textbf{D1. Canadian cost-of-living channels.} Pick \textbf{two} Canadian-relevant channels and write one paragraph on each, connecting them to your inflation surge period:
\begin{itemize}
\item food and groceries,
\item housing costs (rent, mortgage interest),
\item transportation and gasoline,
\item services inflation,
\item exchange rate and imported goods.
\end{itemize}

\vspace{0.5em}
\textbf{Channel 1 paragraph:}\\[0.75em]
\underline{\hspace{0.9\linewidth}}\\[0.75em]
\underline{\hspace{0.9\linewidth}}\\[0.75em]
\underline{\hspace{0.9\linewidth}}

\vspace{0.5em}
\textbf{Channel 2 paragraph:}\\[0.75em]
\underline{\hspace{0.9\linewidth}}\\[0.75em]
\underline{\hspace{0.9\linewidth}}\\[0.75em]
\underline{\hspace{0.9\linewidth}}

\section*{Part E. Optional extension (5 minutes)}
Compute an approximate \textbf{real policy rate}:
\[
r_t \approx i_t - \pi_t,
\]
where $i_t$ is the policy rate and $\pi_t$ is year-over-year inflation.

\vspace{0.5em}
\textbf{E1.} Identify one period where the approximate real policy rate is negative. In Canada, who is more likely to benefit from negative real rates (borrowers or savers), and why?\\[0.75em]
\underline{\hspace{0.9\linewidth}}\\[0.75em]
\underline{\hspace{0.9\linewidth}}\\[0.75em]
\underline{\hspace{0.9\linewidth}}

\section*{2-minute group share-out (if time permits)}
Each group shares: (i) the month of peak inflation, (ii) whether policy appears to lag or lead, and (iii) one Canadian cost-of-living channel that is most salient in your charts.

\end{document}
